% !TEX root = master.tex
% !TEX program = lualatex

%%%%%%%%%%%%%%%%%%%%%%%%%%%%%%%%%%%%%%%%%%%%%%%%%%%%%%
\section{DNS}

\subsection{IP Addresses}

An IP address identifies a \textbf{network interface}, not directly a computer.

Each end-system has at least one network interface.

%%%%%%%%%%%%%%%%%%%%%%%%%%%%%%%%%%%%%%%%%%%%%%%%%%%%%%
\subsection{Port Numbers}

Each process has a unique port number on its host.

Examples:

\begin{itemize}
\item 80 → HTTP
\item 443 → HTTPS
\end{itemize}

These are called \textbf{reserved ports}.

%%%%%%%%%%%%%%%%%%%%%%%%%%%%%%%%%%%%%%%%%%%%%%%%%%%%%%
\subsection{URL Structure}

Example:

\begin{verbatim}
https://www.youtube.com/watch?v=...
\end{verbatim}

Components:

\begin{itemize}
\item protocol
\item hostname
\item path
\end{itemize}

%%%%%%%%%%%%%%%%%%%%%%%%%%%%%%%%%%%%%%%%%%%%%%%%%%%%%%
\subsection{Domain Name System (DNS)}

DNS translates:

\[
\text{hostname} \leftrightarrow \text{IP address}
\]

Example:

\[
www.youtube.com \rightarrow 192.178.170.91
\]

DNS is itself a distributed application.

%%%%%%%%%%%%%%%%%%%%%%%%%%%%%%%%%%%%%%%%%%%%%%%%%%%%%%
\subsection{DNS Client and Server}

Components:

\begin{itemize}
\item DNS client (resolver)
\item DNS server
\end{itemize}

Server port:

\[
53
\]

%%%%%%%%%%%%%%%%%%%%%%%%%%%%%%%%%%%%%%%%%%%%%%%%%%%%%%
\subsection{DNS Hierarchy}

DNS uses a hierarchical architecture.

Levels:

\begin{itemize}
\item Root servers
\item Top-Level Domain (TLD) servers
\item Authoritative DNS servers
\end{itemize}

Example resolution:

\begin{enumerate}
\item Client asks local DNS server
\item Local server asks root server
\item Root server refers to TLD server
\item TLD server refers to authoritative server
\item Authoritative server returns IP address
\end{enumerate}


